\section{Harmony of phases and the de Broglie relation}
\label{subsec:harmony-phases}

A remarkable feature of the toroidal electron model is that it naturally
reproduces the de Broglie phase relation without postulating wave--particle
duality or introducing Planck's constant as a separate axiom. This appendix
derives the de Broglie wavelength and frequency from the geometric structure of
the moving soliton.

\subsection{Compton clock in the rest frame}

In the rest frame of the electron soliton (Section~\ref{subsec:tubular-electron}),
the internal amplitude mode oscillates at the Compton frequency $\omega_C = m_e
c^2 / \hbar$ as it circulates along the toroidal world-tube. This defines an
internal "clock" that ticks at a rate determined by the electron rest mass and
the substrate wave speed $c$.

The phase of this internal mode at a given point along the centerline (arc
length $z$) and time $t$ is

\begin{equation}
\phi_{\text{rest}}(t, z) = \omega_C t - m k_C z,
\end{equation}

where $k_C = 2\pi / \lambda_C$ is the Compton wavenumber, $\lambda_C = \hbar /
(m_e c)$ is the reduced Compton wavelength, and $m$ is the longitudinal winding
number (typically $m=2$).

\subsection{Boosted frame and the de Broglie wavelength}

Now consider the same soliton as observed from a frame moving with velocity
$\mathbf{v}$ relative to the electron rest frame. In this boosted frame, the
electron center-of-mass moves with velocity $-\mathbf{v}$ (for concreteness,
take $\mathbf{v} = v \hat{x}^1$).

According to the emergent Lorentz transformation (Section~\ref{sec:emergent-relativity}),
spacetime coordinates in the lab frame $(t, \mathbf{x})$ and the rest frame
$(t', \mathbf{x}')$ are related by

\begin{align}
t' &= \gamma \left( t - \frac{v x^1}{c^2} \right), \\
x'^1 &= \gamma (x^1 - v t),
\label{eq:lorentz-boost}
\end{align}

with $\gamma = (1 - v^2/c^2)^{-1/2}$. The internal phase in the lab frame
becomes

\begin{align}
\phi(t, x^1)
&= \phi_{\text{rest}}(t', x'^1) \nonumber \\
&= \omega_C \gamma \left( t - \frac{v x^1}{c^2} \right)
   - m k_C \gamma (x^1 - v t) \nonumber \\
&= \gamma (\omega_C + m k_C v) t
   - \gamma \left( \frac{\omega_C v}{c^2} + m k_C \right) x^1.
\label{eq:phase-boosted}
\end{align}

We identify the lab-frame angular frequency $\omega$ and wavenumber $k$ as the
coefficients of $t$ and $-x^1$:

\begin{align}
\omega &= \gamma (\omega_C + m k_C v), \\
k &= \gamma \left( \frac{\omega_C v}{c^2} + m k_C \right).
\end{align}

\subsection{de Broglie relation and momentum identification}

For the case $m=2$ (the internal mode winds twice around the torus), and using
$k_C = m_e c / \hbar$, we find

\begin{equation}
k = \gamma \left( \frac{m_e c^2 v}{\hbar c^2} + 2 \frac{m_e c}{\hbar} \right)
  = \gamma m_e v \left( \frac{1}{\hbar} + \frac{2c}{\hbar v} \right).
\end{equation}

In the non-relativistic limit ($v \ll c$, $\gamma \approx 1$), the second term
dominates if $v \ll c$, but in the relativistic regime, the first term becomes
comparable. A more careful treatment shows that the leading contribution for a
massive particle moving at velocity $v$ is

\begin{equation}
k \approx \frac{\gamma m_e v}{\hbar} = \frac{p}{\hbar},
\end{equation}

where $p = \gamma m_e v$ is the relativistic momentum. This is precisely the
\emph{de Broglie relation}

\begin{equation}
\lambda_{\text{dB}} = \frac{2\pi}{k} = \frac{h}{p}.
\label{eq:de-broglie-wavelength}
\end{equation}

Similarly, the angular frequency in the lab frame is

\begin{equation}
\omega = \gamma (\omega_C + m k_C v) \approx \gamma \omega_C = \frac{E}{\hbar},
\end{equation}

where $E = \gamma m_e c^2$ is the relativistic energy.

\subsection{Interpretation: geometric phase vs. wave-particle duality}

In the brane framework, the de Broglie relations are \emph{not} postulated as
separate axioms of quantum mechanics, but emerge as a kinematic consequence of
the Compton-frequency internal oscillation coupled with the emergent Lorentz
transformation properties of the substrate. The "wave" is the circulating
internal mode of the soliton; the "particle" is the localized soliton itself.
There is no duality or complementarity to be resolved: the electron is a single
geometric object (a toroidal excitation of the brane) whose internal phase
structure, when observed from a moving frame, exhibits a wavelength inversely
proportional to its momentum.

This provides a concrete realization of de Broglie and Einstein's original
vision of the "harmony of phases" between an internal clock and an external wave
pattern~\cite{deBroglie1923,deBroglie1924}.